% =====================================================================
% Capítulo 6. Resultados y discusión               Presupuesto: 14 páginas
% ---------------------------------------------------------------------
% Cada sección sigue el mismo patrón (párrafo por paso):
%   1. Qué pregunta responde el experimento (una oración).
%   2. La figura o tabla, citada en el texto ANTES de que aparezca.
%   3. Qué se ve, con los números clave (no repetir toda la tabla).
%   4. Qué significa y qué NO se puede afirmar.
% Los resultados negativos también se reportan.
% Cada número debe salir de un CSV del repo; anota la fuente en el comentario
% de la figura.
% =====================================================================
\chapter{Resultados y discusión}
\label{cap:resultados}

\section{Elección del diseño (2022)}
\label{sec:res-escalera}

\section{Elección del punto de operación}
\label{sec:res-punto}

\section{Una estación atacada: NERAD contra PARAAD-SA}
\label{sec:res-una}
% Figura de prueba de la paleta (borrar o conservar cuando lleguemos aquí):
\begin{figure}[htbp]
  \centering
  % figuras/recall_banco_p99.tex
% Recall con una estación atacada (escenario E1), umbral p99 de 2023.
% Fuente: results/linea_base/rev24/metricas_banco.csv, filas E1, nivel 0.01,
%   detectores Ens-Vb-X0.2596 (NERAD) y PARAAD. Verificado el 2026-11-27.
\begin{tikzpicture}
\begin{axis}[
  ejes tesis,
  width=0.9\textwidth, height=6cm,
  ybar, bar width=9pt,
  ymin=0, ymax=1.05, ytick={0,0.2,...,1},
  ylabel={Recall},
  xlabel={Magnitud de la anomalía},
  symbolic x coords={d20,d50,d100,a10,a20,a40},
  xtick=data,
  xticklabels={$\pm$20\,\%, $\pm$50\,\%, $\pm$100\,\%, $\pm$10 ppb, $\pm$20 ppb, $\pm$40 ppb},
  enlarge x limits=0.1,
  legend style={at={(0.02,0.98)}, anchor=north west},
  legend cell align=left,
]
\addplot[barra=ensX] coordinates {(d20,0.2020) (d50,0.8019) (d100,0.9319) (a10,0.2279) (a20,0.8312) (a40,0.9584)};
\addlegendentry{NERAD}
\addplot[barra=paraad] coordinates {(d20,0.0884) (d50,0.3647) (d100,0.6753) (a10,0.1013) (a20,0.5656) (a40,0.9201)};
\addlegendentry{PARAAD-SA}
\end{axis}
\end{tikzpicture}

  \caption[Recall en el banco de pruebas con umbral p99]{Recall de NERAD y
  PARAAD-SA con una estación atacada en el banco de pruebas (2024 a julio de
  2026, 30 calendarios, 1\,777 estación-días atacados), con el umbral p99
  calibrado en 2023 sin ataques.}
  \label{fig:recall-banco-p99}
\end{figure}

\section{Despliegue en agosto y septiembre de 2026}
\label{sec:res-despliegue}

\section{Ataque coordinado a las vecinas}
\label{sec:res-coordinado}

\section{Sesgo sostenido}
\label{sec:res-sesgo}

\section{Resultados negativos}
\label{sec:res-negativos}
% Examen 2017--2019, EP1, T0, fusión con PARAAD-SA.

\section{Exploración con imágenes TINTO}
\label{sec:res-tinto}

\section{Amenazas a la validez}
\label{sec:validez}
% Anomalías sintéticas; el banco orientó revisiones; despliegue descriptivo;
% reimplementación de PARAAD-SA; ventana de despliegue fuera de la temporada alta.
